% This template was partially made with assistance from Microsoft Co-pilot

% ------------------------------------------------------------------
% CAPSTONE SOFTWARE REQUIREMENTS SPECIFICATION (SRS)

% 
% This template is inspired by IEEE 29148 but simplified for use in 
% a senior software engineering capstone project. 
% 
% The goal is not to maximize document length. 
% The goal is to produce a clear, testable, and traceable set of 
% requirements that can guide design, implementation, and testing. 
% 
% Recommended workflow: 
% Stakeholder Requirements 
% -> 
% Use Cases 
% -> 
% Functional Requirements 
% -> 
% Acceptance Tests (will be in your VnV plan)
% -> 
% Traceability Matrix (will excluded traceability for tests, that will be in the VnV plan)
% 
% Students should update this document throughout the project. 
% ------------------------------------------------------------------

\documentclass[12pt]{article}

\usepackage{amsmath, mathtools}
\usepackage{amsfonts}
\usepackage{amssymb}
\usepackage{graphicx}
\usepackage{colortbl}
\usepackage{xr}
\usepackage{hyperref}
\usepackage{longtable}
\usepackage{xfrac}
\usepackage{tabularx}
\usepackage{float}
\usepackage{siunitx}
\usepackage{booktabs}
\usepackage{caption}
\usepackage{pdflscape}
\usepackage{afterpage}

\usepackage[round]{natbib}

%\usepackage{refcheck}

\hypersetup{
    bookmarks=true,         % show bookmarks bar?
      colorlinks=true,       % false: boxed links; true: colored links
    linkcolor=red,          % color of internal links (change box color with linkbordercolor)
    citecolor=green,        % color of links to bibliography
    filecolor=magenta,      % color of file links
    urlcolor=cyan           % color of external links
}

\input{../Comments.text}
\input{../Common.text}

% Used so that cross-references have a meaningful prefix

\newcounter{srnum} %Stakeholder Requirement (SR) Number
\newcommand{\srthesrnum}{SR\thesrnum}
\newcommand{\srref}[1]{SR\ref{#1}}

\newcounter{ucnum} %Use Case (UC) Number
\newcommand{\uctheucnum}{UC\theucnum}
\newcommand{\ucref}[1]{UC\ref{#1}}

\newcounter{reqnum} %Functional Requirement (FR) Number
\newcommand{\rthereqnum}{FR\thereqnum}
\newcommand{\rref}[1]{FR\ref{#1}}

\newcounter{nfrnum} %Non-Functional Requirement (NFR) Number
\newcommand{\rthenfrnum}{NFR\thenfrnum}
\newcommand{\nfrref}[1]{NFR\ref{#1}}

\newcounter{lcnum} %Likely Change (LC) Number
\newcommand{\lthelcnum}{LC\thelcnum}
\newcommand{\lcref}[1]{LC\ref{#1}}

\usepackage{fullpage}

\begin{document}

\title{Software Requirements Specification for \progname: subtitle describing software} 
\author{\authname}
\date{\today}
	
\maketitle

~\newpage

\pagenumbering{roman}

\tableofcontents

~\newpage

\section*{Revision History}

\begin{tabularx}{\textwidth}{p{3cm}p{2cm}X}
\toprule {\bf Date} & {\bf Version} & {\bf Notes}\\
\midrule
Date 1 & 1.0 & Notes\\
Date 2 & 1.1 & Notes\\
\bottomrule
\end{tabularx}

~\\
\plt{All capstone SRS documents should have a reflection appendix.}

\newpage

\pagenumbering{arabic}

\plt{This SRS template is based on the IEEE template, with input from co-pilot
  and S. Smith. It will get you started.  You should not modify the section
  headings, without first discussing the change with the course instructor.
  Modification means you are not following the template, which loses some of the
  advantage of a template, especially standardization.}

\plt{Feel free to change the appearance of the report by modifying the LaTeX
  commands.}

\plt{This template document assumes that a single program is being documented.
  If you are documenting a family of models, you should start with a commonality
  analysis.}

\plt{The SRS is not generally written, or read, sequentially.  The SRS is a
  reference document.  It is generally read in an ad hoc order, as the need
  arises.}

\plt{Guiding principles for the SRS document:
\begin{itemize}
\item Do not repeat the same information at the same abstraction level.  If
  information is repeated, the repetition should be at a different abstraction
  level.
\end{itemize}
}

\plt{The template description comments should be disabled before submitting this
  document for grading.}

\plt{You can borrow any wording from the text given in the template.  It is part
  of the template, and not considered an instance of academic integrity.  Of
  course, you need to cite the source of the template.}

\plt{When the documentation is done, it should be possible to trace back to the
  source of every piece of information.  Some information will come from
  external sources, like terminology.}

\plt{An SRS document should have the following qualities: unambiguous,
  consistent, complete, validatable, abstract and traceable.}

\plt{The overall goal of the SRS is that someone can learn,
  understand and verify the captured domain knowledge.  They should not have to
  trust the authors of the SRS on any statements.  They should be able to
  independently verify/derive every statement made.}

% ============================================================== 
\section{Introduction} 
% ==============================================================

\plt{Explain what system is being specified. State purpose, users, and project goals.}

\plt{Introduce the name of your software as \progname{}.  Using \progname{} as
a command lets you change the name if necessary and have the change propagate to
all documents.}

\plt{Include a reference to the template this work is based on.  \citet{ISOIECIEEE29148_2018}
justifies the SRS structure and \citet{ISOIECIEEE15288_2023} provides the
broader life-cycle context that underlines modern requirements engineering
practice.}

\subsection{Purpose}

\plt{State the purpose of the system.}

\plt{The purpose should be in the context of the project itself, not in the
  context of this course. Although the ``purpose'' of the document is to get a
  grade, you should not mention this.  Instead, ``fake it'' as if this is a real
  project.}

\progname{} is intended to solve ... \plt{What problem does your program solve?
The description here should be in the problem space, not the solution space.}

\subsection{Scope}
\plt{Briefly describe:
\begin{itemize}
  \item Problem being solved
  \item Intended users
  \item Major functionality
  \item Expected benefits
  \item Clarifying the scope often means saying what the software doesn't do
\end{itemize}
}

\plt{The scope section is relevant for understanding the boundaries of the system. It is also relevant for later determining typical values of inputs. The scope should make it clear what inputs are reasonable to expect. This is a distinction between scope and context (context is a later section).  Scope affects the inputs while context affects how the software will be used.}

TODO

\subsection{Stakeholders}
\begin{itemize}
\item Customer
\item End Users
\item Project Sponsor
\item Development Team
\item Administrators
\end{itemize}

\subsection{Definitions and Acronyms}

\plt{This section provides the meaning of the different words and phrases used
  in the domain of the problem.}

\begin{longtable}{p{0.25\textwidth} p{0.65\textwidth}}
\toprule
Term & Definition\\
\midrule
TODO & TODO\\
\bottomrule
\caption{Definitions and Acronyms} \label{tab:DefinitionsAndAcronyms}
\end{longtable}

\subsection{Characteristics of Intended Reader} \label{sec_IntendedReader}

\plt{This section summarizes the skills and knowledge of the readers of the
  SRS.  It does NOT have the same purpose as the ``User Characteristics''
  section (Section~\ref{SecUserCharacteristics}).  The intended readers are the
  people that will read, review and maintain the SRS.  They are the people that
  will conceivably design the software that is intended to meet the
  requirements.  The user, on the other hand, is the person that uses the
  software that is built.  They may never read this SRS document.  Of course,
  the same person could be a ``user'' and an ``intended reader.''}

\plt{The intended reader characteristics should be written as unambiguously and
  as specifically as possible.  Rather than say, the user should have an
  understanding of physics, say what kind of physics and at what level.  For
  instance, is high school physics adequate, or should the reader have had a
  graduate course on advanced quantum mechanics?}

% ============================================================== 
\section{System Overview}

This section provides general information about the system.  It identifies the
problem, the vision, and the interfaces between the system and its environment.
This section also describes user characteristics and lists the system
constraints.

\plt{The purpose of this section is to provide general information about the
  system so that the specific requirements in the later sections will be easier to
  understand. The general system description section is designed to be
  changeable independent of changes to the functional requirements documented in
  the specific system description. The general system description provides a
  context for a family of related models.  The general description can stay the
  same, while specific details are changed between family members.}

\subsection{Problem Statement}

\plt{Current situation, problem, impact, desired improvement.  You will likely borrow some text from your problem description document.}

TODO

\subsection{Product Vision}

\plt{A clear and concise statement of the product's purpose and goals.}

\plt{The goal statements are functional and nonfunctional objectives that system
  under consideration should achieve. Goals provide criteria for sufficient
  completeness of a requirements specification and for requirements pertinence.
  Goals will be refined by requirements.  The goals are written abstractly, with
  a minimal amount of technical language.  They should be understandable by
  non-domain experts.}

\plt{You will likely borrow from your goal statements in your problem
description document.}

TODO
  
\subsection{System Context}

\plt{Include users, external systems, APIs, databases, and third-party services.}

\plt{Your system context will include a figure that shows the abstract view of
  the software.  In some cases the diagram will show other external
  entities, besides the user.  For instance, when the software product is a
  library, the user will be another software program, not an actual end user.
  If there are system constraints that the software must work with external
  libraries, these libraries can also be shown on the System Context diagram.
  They should only be named with a specific library name if this is required by
  the system constraint.}

\begin{figure}[h!]
\begin{center}
 \includegraphics[width=0.6\textwidth]{SystemContextFigure}
\caption{System Context}
\label{Fig_SystemContext} 
\end{center}
\end{figure}

\plt{For each of the entities in the system context diagram its responsibilities
  should be listed.  Whenever possible the system should check for data quality,
  but for some cases the user will need to assume that responsibility.  The list
  of responsibilities should be about the inputs and outputs only, and they
  should be abstract.  Details should not be presented here.  However, the
  information should not be so abstract as to just say ``inputs'' and
  ``outputs''.  A summarizing phrase can be used to characterize the inputs.
  For instance, saying ``material properties'' provides some information, but it
  stays away from the detail of listing every required properties.}

\begin{itemize}
\item User Responsibilities:
\begin{itemize}
\item 
\end{itemize}
\item \progname{} Responsibilities:
\begin{itemize}
\item Detect data type mismatch, such as a string of characters instead of a
  floating point number
\item 
\end{itemize}
\end{itemize}

\plt{Identify in what context the software will typically be used.  Is it for
exploration? education? engineering work? scientific work?. Identify whether it
will be used for mission-critical or safety-critical applications.} 

\plt{This additional context information is needed to determine how much effort
should be devoted to the rationale section.  If the application is
safety-critical, the bar is higher.  This is currently less structured, but
analogous to, the idea to the Automotive Safety Integrity Levels (ASILs) that
McSCert uses in their automotive hazard analyses.}

\subsection{User Characteristics} \label{SecUserCharacteristics}

\plt{This section summarizes the knowledge/skills expected of the user.
  Measuring usability, which is often a required non-function requirement,
  requires knowledge of a typical user.  As mentioned above, the user is a
  different role from the ``intended reader,'' as given in
  Section~\ref{sec_IntendedReader}.  As in Section~\ref{sec_IntendedReader}, the
  user characteristics should be specific an unambiguous.  For instance, ``The
  end user of \progname{} should have an understanding of undergraduate Level 1
  Calculus and Physics.''}

\subsection{System Constraints}

\plt{System constraints differ from other type of requirements because they
  limit the developers' options in the system design and they identify how the
  eventual system must fit into the world. This is the only place in the SRS
  where design decisions can be specified.  That is, the quality requirement for
  abstraction is relaxed here.  However, system constraints should only be
  included if they are truly required.}

\plt{This is where you would mention platform constrains, and licensing
restrictions.}

% ============================================================== 
\section{Stakeholder Requirements}

\plt{Technology-independent requirements.  These come from statements made by
the stakeholders.  You will find these through interviews and surveys.  They are
often expressed in natural language.  They are not necessarily testable, but
they should be verifiable.}

\plt{The stakeholder requirements are numbered so that you can have traceability to other parts of your document.}

\noindent \begin{itemize}

%\item[SR\refstepcounter{srnum}\thesrnum \label{OctalChecker}] This program shall check to make sure that the user enters in the correct complimentary octal number. 

%\item[SR\refstepcounter{srnum}\thesrnum \label{SelectionChecker}] This program shall check to ensure that the selected square matches the requirements displayed on the screen (ex. if it asks for the complimentary of a highlighted position, the right one should be selected).

%\item[SR\refstepcounter{srnum}\thesrnum \label{CubeDatabase}] This program shall extract validated octal magic cubes from its own database.

\item[SR\refstepcounter{srnum}\thesrnum \label{MagicCubeValidator}] This program shall validate the stakeholder's answers in terms of both position and octal answer.

\item[SR\refstepcounter{srnum}\thesrnum \label{CubeVisualizer}] This program shall visualize the current octal magic cube being solved by the stakeholder.

\item[SR\refstepcounter{srnum}\thesrnum \label{OctalExplainer}] This program shall teach stakeholders how octal number works.

\item[SR\refstepcounter{srnum}\thesrnum \label{MagicCubeExplainer}] This program shall teach stakeholders how magic squares and then cubes work.

\item[SR\refstepcounter{srnum}\thesrnum \label{CompletionVerifier}] This program shall verify if the stakeholders satisfied all the requirements of a level.

\item[SR\refstepcounter{srnum}\thesrnum \label{LevelGenerator}] This program shall show stakeholders different cubes with tasks of the same difficulty at the same level.

\item[SR\refstepcounter{srnum}\thesrnum \label{ScoreKeeper}] This program shall time stakeholders and show them their current and final score.


\end{itemize}

\plt{You should give your SRs meaningful labels, so that you can refer to them
in other parts of the SRS.  The label names should be meaningful mnemonics, not
numbers.}

% ============================================================== 
\section{User Personas}
\plt{Create 2-5 realistic personas.
  Personas are fictional characters that represent the different user types that
  might use your software.  They are used to help understand the users'
  needs, experiences, behaviors and goals.  They are not real people, but they
  should be based on real data collected from user research.}

\subsection{Persona 1}
\begin{itemize}
\item Name: Scotty Thompson
\item Age: 16
\item Occupation: Student
\item Location: Toronto, Canada
\item Description: Scotty is a Grade 10 student who stumbled upon this program when he was learning octal numbers in his Computer Technology class. His teacher took a glance at the program and recommended using it to study for the octal/hexadecimal test. The only understanding he has of the octal system is from his computer technology course, which he pays attentions to. but is lost in his train of thoughts when learning about concepts that don't apply to everyday life.
\item Goals: 
\begin{itemize}
  \item With a basic undertanding of octal numbers, Scotty wishes to learn how to do them faster (like a multiplication table), so he doesn't get stuck on test questions requiring addition/subtraction.
  \item He doesn't want a simple math checker, he wants more information on why he got the answer wrong, alongside a visual demonstration on how to get the correct answer.
\end{itemize}
\item Frustrations:
\begin{itemize}
  \item He doesn't care about the "magic" and intrinsicality of the magic cube, as he feels like he has no use of learning about these patterns.
  \item He will get bogged down if shown literature about magic cubes as he is using this program for only one goal.
\end{itemize}
\item Technical Ability: He is a gamer who has played browser games since he was a child and is able to quickly click and type.
\end{itemize}

\subsection{Persona 2}
\begin{itemize}
\item Name: John Fye
\item Age: 73
\item Occupation: Retired (Former Stenographer)
\item Location: Hamilton, Canada
\item Description: John is a former court stenographer who recently got a kick out of learning about John Hendricks's work. He wishes to test out his ability to see the patterns in magic cubes as he has been doing base-10 magic square puzzles in his puzzle book for the last 20 years.
\item Goals:
\begin{itemize}
  \item He wants to learn more about octal numbers, as he has a very basic understanding about it, but a lot of Mr. Hendricks's work finds patterns in different bases, and he feels like learning about octal will help him appreciate Mr. Hendricks's accomplishements, such as order 16 perfect magic tesseract.
  \item He wants to discover the patterns in the cube, such as how you can use a starting number to find other numbers (such as the complimentary).
  \item He really wants to dive deep into 
\end{itemize}
\item Frustrations:
\begin{itemize}
  \item The only thing he currently knows about octal is how it cycles from 0 to 7, he never did any math with it.
  \item He struggles to look at a screen for a long time and is usually zoomed in. He takes extensive breaks, which might lead him to forgetting where he left off and what was his train of thoughts at that point.
\end{itemize}
\item Technical Ability: He was a court stenographer, which allows him to type fast. However, he never really used a computer before, which will be a problem when he is using a mouse. He also doesn't want to look at a mostly white screen when using the program.
\end{itemize}


\subsection{Persona 3}
\begin{itemize}
\item Name: Charles Dean
\item Age: 23
\item Occupation: Masters in English
\item Location: Halifax, Canada
\item Description: Charles is someone that starts off their day playing the "Dles", which is a collection of single-level browser games that refreshes daily, like Wordle. He is tired of the flood of general knowledge trivia games, and wants to play a game that tests his brain without needing to know stuff beforehand.
\item Goals:
\begin{itemize}
  \item He wants a daily game that he can complete as part of his daily rotation of games (shouldn't be too long).
  \item He wants the game to run on his browser with minimal load times, not really interested in getting stuck due to the 3D visualization.
\end{itemize}
\item Frustrations:
\begin{itemize}
  \item He isn't really interested in math, so might lose interest if the puzzles are too complex.
  \item He needs a hook (either visually or with flair) that ensures that he plays it more than once as he doesn't start off interested in the subject matter.
\end{itemize}
\item Technical Ability: He is very familiar with browser games and while he hasn't done math since high school, he does remember an inkling of what hexadecimal is from his high school computer course. 
\end{itemize}


% ============================================================== 
\section{Use Cases}
\plt{Use cases are a way to describe the interactions between users and the system.
  They are typically written in natural language and describe the steps that a
  user takes to achieve a specific goal.  Use cases can be used to identify
  functional requirements and to create test cases.}  
\plt{Typically 8-15 substantial use cases.}

\refstepcounter{ucnum} \label{UC_Label1}
\subsection{\ucref{UC_Label1}: Round Setup}

\textbf{Primary Actor}

User

\textbf{Preconditions}
\begin{itemize}
\item Browser game is deployed on an accessible website
\item Game has access to database containing cube configurations
\item No rounds is currently in progress
\end{itemize}

\textbf{Main Success Scenario}
\begin{enumerate}
\item User presses START
\item System wipes screen clean and pings cube database wih appropiate query
\item Cube database sends the values of the cells alongside with a cube number (as an identifier)
\item System creates the 3D visualization and discretely populates visualization cells with its correct, hidden values
\item System prompts user to choose the scoring system (decimal or octal)
\item User clicks on the scoring system they want and its confirmation
\item System selects the appropiate scoring system
\item System randomly populates a single visualization cell with its correct octal and has a final popup describing the next level
\end{enumerate}

\textbf{Alternative Flows}
\begin{itemize}
\item If Cube database doesn't ping back, the System will already have one default cube formation stored locally in its code, so that cube formation will be used instead. 
The System will continuing pinging the database every 30 seconds to see if the database is active and notifiy its developers the first time it doesn't get an answer.
\item If the 3D visualization fails, either fall back to a simpler visualization method or show an error.
\item If the scoring system fails, continue with the steps after displaying a message about how scores isn't going to be tracked for this round.
\end{itemize}

\textbf{Postconditions}
\begin{itemize}
\item All errors and log-type messages will appear as notifications starting from the bottom on a corner.
\item There is a black 3D visualization of a cube in the centre.
\item There is a notification showing '512 of 512 cells loaded' and there is a single octal number on the cube.
\item On one of the upper corners, there is a scoring component consisting of an unstarted timer and a score integer.
\end{itemize}

\refstepcounter{ucnum} \label{UC_Label2}
\subsection{\ucref{UC_Label2}: Tasks and Answer Checking}

\textbf{Primary Actor}

User

\textbf{Preconditions}
\begin{itemize}
\item The system should be set up according to the postconditions of \ucref{UC_Label1}
\end{itemize}

\textbf{Main Success Scenario}
\begin{enumerate}
\item User chooses to begin the levels by clicking START.
\item System displays a textbox asking the user to select the complement of the single populated cell
\item User moves the 3D visualization using the chosen movement method in order to locate the complement
\item User clicks on the cell representing the complement, thereby selecting it
\item System flashes white and displays the properties of the intial populated cell and a textbox for the answer
\item User enters in the complement's value in the Addition/Subtraction Box
\item System alerts the scoring syetem by sending a query with the level of the task that was just finished
\item System shows the user a level overview with a difficulty slider
\item User chooses difficulty
\item System populates the cells according to the difficulty selected and displas a START button
\item User clicks START and selects the cells that corresponds to the half-orthogonal for the starting cell
\item User clicks the DONE button
\item System displays the Addition/Subtraction Box for selected cells
\item User enters in the octal values and clcks DONE
\item System queries the scoring system with information of the level just completed
\item Previous steps repeat for the remaining steps: orthogonal, diagonal on continuous plane, diagonal not on continuous plane, another diagonal not on continuous plane, triagonal, split triagonal 
\item System asks scoring sytem for final score
\item Scoring sstem calculates final score
\item System ends after displaying a final score and buttons to navigate to other sections
\end{enumerate}

\textbf{Alternative Flows}
\begin{itemize}
\item User is able to pull up level overview when they are in the level by selecting an arrow poking out of the side if they are stuck
\item If the User clicks DONE after selecting the wrong cells, the game will flash red and deselect the selected squares, and siplays a message saying they selected the wrongs squares
\item When the user inputs the wrong answer in the Addition/Subtracton box, the box will flash red and clear the incorrect answer, letting them try again. The System will also notify the scoring system that an incorrect answer has been given. 
\item If the 3D visualization is stalling (e. due to the user running out of memory), the System shows a toggle to switch to the 2D view of the cube
\item If the user is signed in, the are able to save the score if they want to
\end{itemize}

\textbf{Postconditions}
\begin{itemize}
\item System should display the final score front and centre
\item There should be an option to either create a new account (while being able to save the current score), sign in, or view account information 
\item There are options displayed to either start a new run, go to the tutorial, or enter other modes
\end{itemize}


\refstepcounter{ucnum} \label{UC_Label3}
\subsection{\ucref{UC_Label3}: Scoring and controls}

\textbf{Primary Actor}

System

\textbf{Preconditions}
\begin{itemize}
\item The system is currently active and in a viable state of execution
\item The scoring system is an instance which is volatile and doesn't keep track of information
\item The scoring system is well integrated with near-instantaneous signals from the system
\item The scoring system is able to use other APIs to ensure that its time signals is correct
\end{itemize}

\textbf{Main Success Scenario}
\begin{enumerate}
\item System queries the scoring system with level number
\item Scoring system initializes a stopwatch, error counter and correct scoring function
\item System pings the scoring system, telling them that the level has started
\item Scoring system starts the timer and listens for signals indicating errors or completion 
\item System pings for an error
\item Scoring system increments error counter
\item System pings for the completion of the level
\item Scoring system runs appropiate scoring function with time and error counter as inputs
\item Scoring system returns the score and its instance terminates
\end{enumerate}

\textbf{Alternative Flows}
\begin{itemize}
\item If the user selected a different scoring base or level, the system (not scoring system) will add the corresponding values for each addition:
\begin{center}
\begin{tabular}{||c c c||} 
 \hline
 Level & Decimal & Octal \\ [0.5ex] 
 \hline\hline
 1 & 10 & 14 \\ 
 \hline
 2 & 20 & 30 \\
 \hline
 3 and 4 & 30 & 44 \\
 \hline
 5 and 6 & 40 & 60 \\
 \hline
 7 and 8 & 60 & 110 \\ [1ex] 
 \hline
\end{tabular}
\end{center}
\item If the user is able to enter in additional numbers, they are rewarded 10 points for desimal and 14 points for octal.
\item If the user clicks STOP/EXIT, the scoring system instance is terminated without any return
\item If the user clicks PAUSE, the stopwatch in the scoring system will pause and keep track of the time paused (so it doesn't mess up when checked with eternal APIs), but the scoring system instance will continue running
\end{itemize}

\textbf{Postconditions}
\begin{itemize}
\item The score given can be verified as base score * (1350/seconds) + score for correct cells/octals for an individual level
\item The system's internal final score is incrememented between levels
\item At the end, the system shall show the internal final score (verified as last internal score + level 8 score)
\end{itemize}

\refstepcounter{ucnum} \label{UC_Label4}
\subsection{\ucref{UC_Label4}: Saving and Recovery}

\textbf{Primary Actor}

User

\textbf{Preconditions}
\begin{itemize}
\item System should be in a valid state of execution
\item User database should be active (in the sense that it exists and connectable) and triggerable
\item The connection between the database and system should be closed.
\end{itemize}

\textbf{Main Success Scenario}
\begin{enumerate}
\item User selects cells and pressed DONE
\item System saves which cells were selected
\item System verifies that the correct cells were selected
\item System saves that the correct cells were selected
\item System calculated how much points the user gets for the selection 
\item System saves the number of points
\item System displays the Addition/Subtraction box based on saved game state
\item User completes level
\item System caculates and saves final score
\item User chooses to save the score
\item System pings user database with user information
\item User database pings back (saying that it is active)
\item System shares final score with user database
\item User database saves score
\item System closes the connection with the user database
\item User clicks to return to main menu 
\item System refreshes its information via user database
\item System displays the main menu with user information (including latest/highest final score) 
\end{enumerate}

\textbf{Alternative Flows}
\begin{itemize}
\item If a save fails (either internal or with a database), the system will retry 3 times. If the system still fails to save, it will notify the user and warn that any further progress might not be saved. The user is free to keep playing and the system will retry saving either every 30 seconds or when the user reaches the net save point (whichever comes first).
\item If the program crashes, after exhausting all possible recovery options, the system will save the last recoverable state in the browser cache. It will then recommend the user refresh or close the browser page. The next time the user opens the page, it will ask the user if it wants to load the previous progress or start anew.
\end{itemize}

\textbf{Postconditions}
\begin{itemize}
\item The connection between the user database and system should be closed
\item The schema and integrity of the user database should be equivalent to its precondition 
\item The system shouldn't have any saved data about internal levels (such as level scores and cell selection) after finishing all the levels
\item The browser cache would optimall always be empty unless there has been a crash and the user hasn't returned to the game
\end{itemize}


\refstepcounter{ucnum} \label{UC_Label5}
\subsection{\ucref{UC_Label5}: 3D Visualization}

\textbf{Primary Actor}

System

\textbf{Preconditions}
\begin{itemize}
\item System is in a valid state of execution
\item 3D visualization is a plugin (removing it just causes the system to use another visualization) 
\item 3D visualization doesn't have any data inputted into it (it starts anew everytime a cube is selected)
\end{itemize}

\textbf{Main Success Scenario}
\begin{enumerate}
\item System gets the cube values and number from the cube database
\item 3D visualization starts to initialize in the middle of the screen
\item System chooses a starting position
\item 3D visualization filled with the correct octal number at that position
\item System repositions the cube to the centre of the screen
\item 3D visualization centers the model at the filled cell
\item User moves mouse across a face of the cube
\item 3D visualization rotates accordingly across its face due to mouse
\item User types in a coordinate in a navigator bar
\item System repositions the cube to the centre of the screen
\item 3D visualization centers the model at that coordinate
\item User presses down on the "s" or "down" key on the keyboard
\item 3D visualization rotates down accordingly across its face due to keyboard
\item User drags mouse on an area off the cube
\item System repositions towards the mouse (so that the cube is off centre)
\item User scrolls down on the mousewheel
\item 3D visualization moves along the z-axis, so that the second and next layers of the cube is exposed
\item User selects a cell
\item System saves which cell has been selected and checks its correctness
\item System tells the 3D visualization which cell has been selected and what colour it should be
\item 3D visualization keeps the cell and previous cell as a bright white light source
\item User enters in number in Addition/Subtraction box
\item System checks correctness of the octals inputted
\item System tells the 3D visualization what to project onto the cell's surface 
\item 3D visualization shows the correct octals the user inputted
\item System calculates level score
\item 3D visualization keeps filled cells
\end{enumerate}

\textbf{Alternative Flows}
\begin{itemize}
\item If the user types in a coordinate that would belong on another layer, the system will reposition the cube to the centre of the screen and the 3D visualization centres on the selected xy coordinates. It will then move through each layer until it arrives at the correct z value.
\item If the user is selecting the wrong cell (ex. as part of a triagonal), the system will inform the 3D visualization that the cell chosen is wrong. The 3D visualization will make the cell chosen red.
\item If the user is on a non-face and tries to rotate the cube, the user will be able to see the missing layers either behind or on the R/LHS of the camera, but they will be transparent so the person is aware of where they are located in the cube.
\item If the system is calculating final scores, the 3D visualization isn't needed, so that instance of the 3D visualization plugin is terminated.
\end{itemize}

\textbf{Postconditions}
\begin{itemize}
\item The 3D plugin shouldn't have any data associated with it (such as information about which cells are filled)
\end{itemize}

\refstepcounter{ucnum} \label{UC_Label6}
\subsection{\ucref{UC_Label6}: Guided Learning and Tutorial}

\textbf{Primary Actor}

User

\textbf{Preconditions}
\begin{itemize}
\item The system is in a valid state of execution
\item The system is not in the process of conducting a level/game
\end{itemize}

\textbf{Main Success Scenario}
\begin{enumerate}
\item User selects the tutorial option from the menu
\item System clears screen and projects empty cube
\item System chooses a cube and projects all cells
\item User clicks button to learn about octal
\item System displays popup with an octal explanation, solved examples and calculator
\item User clicks "Test" button
\item System generates random questions and textboxes for solutions
\item User enters solutions
\item System verifies answers and closes popup
\item User clicks button to learn about magic squares
\item System removes navigation controls and centres on a face (thereby a magic cube)
\item System popups explanation about magic square
\item User clicks "Test" button
\item System shows a stored magic square
\item User enters solutions
\item System verifies and shows a harder magic square
\item User exits back to main menu 
\item User starts first game
\item System sets up game
\item User clicks "Hint" button
\item System popups coordinate hint about where the next cell is located
\item User clicks "Hint" button
\item System traces path to the correct square with explanations
\item User selects correct cells
\item System displays the Additon/Subtraction Box
\item User clicks "Hint" button
\item System displays the octal steps to arrive at the correct answer
\end{enumerate}

\textbf{Alternative Flows}
\begin{itemize}
\item If the user asks for more magic square examples, the system displays a link to an educational website.
\item If the user taps the stopwatch, the system will display a message allowing the use to disable the scoring system. If the user wants to, it will result in no scores being recorded.
\item If the user runs out the level timer (22.5 minutes), the system will display a messgae saying that the score wouldn't be counted. The system will not record the score and allow unlimited mistakes and hints.
\end{itemize}

\textbf{Postconditions}
\begin{itemize}
\item The system is in a valid state of execution
\end{itemize}




\refstepcounter{ucnum} \label{UC_Label7}
\subsection{\ucref{UC_Label7}: Create Account}

\textbf{Primary Actor}

User

\textbf{Preconditions}
\begin{itemize}
\item The system is in a valid state of execution
\item The system is not in the process of conducting a level (can sign in between levels)
\item User database should be active (in the sense that it exists and connectable) and triggerable
\item The connection between the database and system should be closed.
\end{itemize}

\textbf{Main Success Scenario}
\begin{enumerate}
\item User clicks button to sign in
\item System saves current game state in browser cache and displays information to sign in
\item User selects "Create Account"
\item System displays a page with a username textbox and username guide
\item User types in appropiate username
\item System verifies username follows username requirements and pings user database
\item User database verifies username doesn't exist
\item Systems displays a page with a password textbox and password guide
\item System verifies password follows password requirements and creates a new user profile and inserts it into the user database
\item System asks user to log in
\item User logs in
\item User database confirms user
\item System recreates game state via browser cache and deletes cache
\item User completes game
\item System saves final score, time needed and date by appending it in user database
\item User database saves and closes connection
\item User goes to main menu
\item System opens connection to user database and displays information
\item User closes browser tab
\item System closes connection to user database
\end{enumerate}

\textbf{Alternative Flows}
\begin{itemize}
\item If the user writes a username or password that doesn't follow the requirements, the system will display an error message alongside an option to type in another name.
\item If a user database ping or save fails, the system will retry 3 times. If the system still fails to get a response to the ping, it will notify the user and say that they failed in making an account. It will also warn that any further progress might not be saved. When the user finishes the final level, the system will retry pinging the user database and if it responds, the system will display a message saying that the user creation page is working again.
\end{itemize}

\textbf{Postconditions}
\begin{itemize}
\item The connection between the user database and system should be closed
\item The schema and integrity of the user database should be equivalent to its precondition 
\end{itemize}


\refstepcounter{ucnum} \label{UC_Label8}
\subsection{\ucref{UC_Label8}: Daily Challenge}

\textbf{Primary Actor}

User

\textbf{Preconditions}
\begin{itemize}
\item The system is in a valid state of execution
\item The system is not in the process of conducting a level/game
\item User and Daily Challenge database (DC) should be active (in the sense that it exists and connectable) and triggerable
\item The connection between the databases and system should be closed.
\end{itemize}

\textbf{Main Success Scenario}
\begin{enumerate}
\item User clicks button to begin Daily Challenge
\item System asks if the user wants to log in
\item User logs in
\item System pings DC
\item DC replies with Daily Challenge  plane colour, orientation, and starter
\item System saves Daily Challenge game state in browser
\item System asks if the user is ready to play
\item User starts challenge
\item System uses cube and starter given instead of pinging cube database and random starter
\item User finishes game
\item System pings user database and DC
\item System saves final score, time needed and date by appending it in DC column in user database
\item System creates a copyable graphic that can be displayed by the user
\item User database saves and closes connection
\item DC sends leaderboard scores to System
\item System checks if the user score belongs on leaderboard
\item System saves final score and time needed at appropiate rank by inserting it into DC's daily leaderboard column
\item System closes connection to DC
\item User goes to main menu
\item System opens connection to user database and displays information
\item User goes to Daily Challenge leaderboard
\item System pings DC for leaderboard data
\item System displays leaderboard with copy button
\item User clicks out
\item System closes DC's connection
\item User closes browser tab
\item System closes connection to user database
\end{enumerate}

\textbf{Alternative Flows}
\begin{itemize}
\item If a DC ping or save fails, the system will retry 3 times. If the system still fails to get a response to the ping, it will notify the user and say that their Daily Challenge attempt will not be counted. Since the user already saw the puzzle, they will not be able to redo the Daily Challenge until it refreshes.
\end{itemize}

\textbf{Postconditions}
\begin{itemize}
\item The connection between the databases and system should be closed
\item The schema and integrity of the bases should be equivalent to its precondition 
\end{itemize}

% ============================================================== 
\section{Functional Requirements}
\plt{Functional requirements describe what the system should do.  They are typically
expressed as a set of conditions or capabilities that the system must have.  These
requirements are usually derived from the stakeholder requirements and use cases.}

\plt{Use SHALL statements. Requirements must be atomic, testable, and traceable.}

\noindent \begin{itemize}

\item[FR\refstepcounter{reqnum}\thereqnum \label{FR_Label1}:] The system shall ...

\item[FR\refstepcounter{reqnum}\thereqnum \label{FR_Label2}:] The system shall ...

\item[FR\refstepcounter{reqnum}\thereqnum \label{FR_Label3}:] The system shall ...

\end{itemize}

\plt{The requirements are given labels so that they can be referenced in other
parts of the SRS.  The label names should be meaningful mnemonics, not numbers.}

% ============================================================== 
\section{Data Requirements}

\plt{The Data Requirements section describes what information the system must manage.
  It includes the major data entities, their attributes, and the relationships
  between them.  It also includes a data dictionary that defines the meaning of
  each attribute.}

\subsection{Major Data Entities}

\plt{What information must exist for the system to accomplish its goals?}

\plt{For each entity, provide a short explanation.}

\begin{itemize}
\item User
\item Song
\item Session
\item Report
\end{itemize}

\subsection{Data Dictionary}

\plt{For each entity, define all important attributes.}

\begin{longtable}{p{0.20\textwidth} p{0.20\textwidth} p{0.45\textwidth}}
\toprule
Attribute & Type & Description\\
\midrule
user\_id & UUID & Unique identifier\\
email & String & User email address\\
\bottomrule
\caption{Data Dictionary} \label{tab:DataDictionary}
\end{longtable}

\subsection{Privacy and Data Retention}

\plt{Discuss storage, retention, deletion, privacy controls,
 consent, encryption, and compliance requirements.}

TODO

% ============================================================== 
\section{Non-Functional Requirements}

\plt{You may have to add other NFRs.}

\plt{List your nonfunctional requirements.  You may consider using a fit
  criterion to make them verifiable.}

\plt{The goal is for the nonfunctional requirements to be unambiguous, abstract
  and verifiable.  This isn't easy to show succinctly, so a good strategy may be
to give a ``high level'' view of the requirement, but allow for the details to
be covered in the Verification and Validation document.}

\plt{An absolute requirement on a quality of the system is rarely needed.  For
  instance, an accuracy of 0.0101 \% is likely fine, even if the requirement is
  for 0.01 \% accuracy.  Therefore, the emphasis will often be more on
  describing now well the quality is achieved, through experimentation, and
  possibly theory, rather than meeting some bar that was defined a priori.}

\plt{You do not need an entry for correctness in your NFRs.  The purpose of the
  SRS is to record the requirements that need to be satisfied for correctness.
  Any statement of correctness would just be redundant. Rather than discuss
  correctness, you can characterize how far away from the correct (true)
  solution you are allowed to be.  This is discussed under accuracy.}

\noindent \begin{itemize}

\item[NFR\refstepcounter{nfrnum}\thenfrnum \label{NFR_Accuracy}:]
  \textbf{Accuracy} \plt{Characterize the accuracy by giving the context/use for
    the software.  Maybe something like, ``The accuracy of the computed
    solutions should meet the level needed for $<$engineering or scientific
    application$>$.  The level of accuracy achieved by \progname{} shall be
    described following the procedure given in Section~X of the Verification and
    Validation Plan.''  If you do this provide a link to the VnV plan.}

\item[NFR\refstepcounter{nfrnum}\thenfrnum \label{NFR_Usability}:] \textbf{Usability}
  \plt{Characterize the usability by giving the context/use for the software.
    You should likely reference the user characteristics section.  The level of
    usability achieved by the software shall be described following the
    procedure given in Section~X of the Verification and Validation Plan.  Include a link to the VnV plan.}

\item[NFR\refstepcounter{nfrnum}\thenfrnum \label{NFR_Maintainability}:]
  \textbf{Maintainability} \plt{The effort required to make any of the likely
    changes listed for \progname{} should be less than FRACTION of the original
    development time.  FRACTION is then a symbolic constant that can be defined
    at the end of the report.}

\item[NFR\refstepcounter{nfrnum}\thenfrnum \label{NFR_Portability}:]
  \textbf{Portability} \plt{This NFR is easier to write than the others.  The
    systems that \progname{} should run on should be listed here.  When possible
    the specific versions of the potential operating environments should be
    given.  To make the NFR verifiable a statement could be made that the tests
    from a given section of the VnV plan can be successfully run on all of the
    possible operating environments.}

\item Other NFRs that might be discussed include verifiability,
  understandability, reusability, performance, reliability and availability, security, and accessibility.  You may have to add other NFRs.

\end{itemize}

% ============================================================== 
\section{AI and Machine Learning Requirements}

\plt{Refer to
\href{https://gitlab.cas.mcmaster.ca/courses/capstone/-/blob/main/Lectures/Lxx_DrMosserGuestLecture/2026_09_Capstone.pdf?ref_type=heads}{Dr.
Mosser's guest lecture} on AI requirements.}

\plt{Include only when applicable.}

\subsection{Model Inputs}

TODO

\subsection{Model Outputs}

TODO

\subsection{Evaluation Metrics}

\plt{Accuracy, Precision, Recall, F1, BLEU, WER, etc.}

TODO

\subsection{Known Limitations}

\plt{Explicitly document expected failure modes.}

TODO

% ============================================================== 
\section{Requirements Traceability Matrices}

\plt{The Requirements Traceability Matrix (RTM) is a table that shows the
relationship between different types of requirements. It is used to ensure that
all requirements are traceable and that no requirements are missed. The typical
flow is: Stakeholder Need -> Use Case -> Requirement}

\plt{There is also a traceability matrix between functional and nonfunctional
  requirements. The purpose of some FRs is to satisfy NFRs.  Therefore, it is
  important to show the relationship between the two.}

\plt{Further traceability to test cases and modules will be given in other
documents.}

\begin{longtable}{p{0.15\textwidth} p{0.15\textwidth} p{0.20\textwidth} p{0.20\textwidth}}

\toprule
Stakeholder Req & Use Case & Functional Req \\
\midrule
\srref{SR_Label1} & \ucref{UC_Label1} & \rref{FR_Label1} \\
\bottomrule

\caption{Traceability Between Stakeholder Requirements, Use Cases, and
Functional Requirements} \label{tab:SR2UC2FRtm}
\end{longtable}

\begin{longtable}{p{0.20\textwidth} p{0.20\textwidth}}

\toprule
Functional Req & Nonfunctional Req \\
\midrule
\rref{FR_Label1} & \nfrref{NFR_Accuracy} \\
\bottomrule
\caption{Traceability Between Functional and Nonfunctional Requirements}
\label{tab:FR2Rtm}

\end{longtable}

% ============================================================== 

\section{Rationale}

\plt{Provide a rationale for the decisions made in the documentation.  Rationale
should be provided for scope decisions, functional requirements, nonfunctioanl
requirements, data requirements, and AI requirements.  You don't need a
rationale for everything, ust the major decisions that are not obvious.  You
should have traceability to the entities in the document that you are
referencing.}

\section{Likely Changes}    

\noindent \begin{itemize}

\item[LC\refstepcounter{lcnum}\thelcnum\label{LC_meaningfulLabel}:] \plt{Give
    the likely changes, with traceability to the entity that is likely to change.}

\end{itemize}

\section{Unlikely Changes}    

\noindent \begin{itemize}

\item[LC\refstepcounter{lcnum}\thelcnum\label{LC_meaningfulLabel}:] \plt{Give
    the unlikely changes.  The design can assume that the changes listed will
    not occur.}

\end{itemize}

\newpage{}

\section{Appendices}

\subsection{Appendix A: Additional Materials}

TODO

\subsection{Appendix B: Development Plan}

\plt{This section is optional.  It is used to explain the plan for developing
  the software.  In particular, this section gives a list of the order in which
  the requirements will be implemented.  In the context of a course  this is
  where you can indicate which requirements will be implemented as part of the
  course, and which will be ``faked'' as future work.  This section can be
  organized as a prioritized list of requirements, or it could show the
  requirements that will be implemented for ``phase 1'', ``phase 2'', etc.}

\subsection{Appendix C: Values of Auxiliary Constants}

\plt{Show the values of the symbolic parameters introduced in the report.}

\plt{The definition of the requirements will likely call for SYMBOLIC\_CONSTANTS.
Their values are defined in this section for easy maintenance.}

\plt{The value of FRACTION, for the Maintainability NFR would be given here.}

TODO

\subsection{Appendix D: Generative AI Disclosure}

\plt{Provide full transparency on how generative AI was used in the creation of
  this document.  This includes any text, figures, or tables that were generated
  by AI.  If any skills were created for AI agents to review the documentation,
  those should be disclosed here.}

\subsection{Appendix E: Reflection}

\plt{This isn't part of the SRS, but it is part of the course.}

The information in this section will be used to evaluate the team members on the
graduate attribute of Lifelong Learning.  

\input{../Reflection.text}

\input{../SRS_Reflection.text}

\newpage

\bibliographystyle {plainnat}
\bibliography {../../refs/References}

\newpage


\end{document}